\section{The coherent contribution in momentum space}
\label{sec:mom}

The results of section~\ref{sec:xs} are written in transverse coordinate space and involve six
transverse integrations. For a part of the cross-section these integrations can be carried out
explicitly. At large $\Nc$ and in the Gaussian approximation \eqref{eq:Gauss}, each quadrupole is the
sum of a factorised term, $S(\x_1,\x_2)S(\x_3,\x_4)$, in which the DA and the CCA are disconnected,
and of a connected term proportional to $F/F$. Keeping only the factorised terms amounts to averaging
the amplitude, rather than the cross-section, over the target field: this is the \emph{coherent}, or
elastic, contribution, in which the target remains in its ground state and the $\qqg$ system is
produced in an overall colour-singlet state, as in coherent diffraction. The connected terms describe
the inelastic contribution and are not considered in this section.

\subsection{Setting up}

For the factorised part, each of the colour structures of section~\ref{sec:colour} becomes a product of
a function of the DA coordinates and a function of the CCA coordinates,
\begin{align}
\cW(\x,\y,\z;\bx,\by,\bz)\big|_{\rm el}&=\big[S(\x,\z)S(\z,\y)-1\big]\big[S(\bz,\bx)S(\by,\bz)-1\big]\,,
\nonumber\\
\cW(\x,\yp,\yp;\bx,\by,\bz)\big|_{\rm el}&=\big[S(\x,\yp)-1\big]\big[S(\bz,\bx)S(\by,\bz)-1\big]\,,
\nonumber\\
\cW(\xp,\y,\xp;\bx,\by,\bz)\big|_{\rm el}&=\big[S(\xp,\y)-1\big]\big[S(\bz,\bx)S(\by,\bz)-1\big]\,,
\label{eq:W-el}
\end{align}
and similarly for the other limits; the structure \eqref{eq:W-int} is already of this form. The
cross-section is then the modulus squared of an amplitude in which the vertex $\cM_1$ ($\tM_1$) is
multiplied by $S(\x,\z)S(\z,\y)-1$ and the vertex $\cM_2$ ($\tM_2$) by $S(\x,\yp)-1$ ($S(\xp,\y)-1$).
We assume that the target is homogeneous in the transverse plane, so that the dipole $S$-matrix depends
only on the dipole size, and introduce its Fourier transform,
\begin{equation}
S(\x,\y)=\int\frac{\dd^2\lv}{(2\pi)^2}\;e^{i\lv\cdot(\x-\y)}\,\cD_S(\lv)\,,\qquad
\cD_S(\lv)=\int\dd^2\bm{r}\;e^{-i\lv\cdot\bm{r}}\,S(\bm{r})\,,\qquad
\int\frac{\dd^2\lv}{(2\pi)^2}\,\cD_S(\lv)=1\,.
\label{eq:DS}
\end{equation}
The momentum $\lv$ is the transverse momentum transferred by the target to the dipole. We use the
shorthand notation
\begin{equation}
\int_{\lv_1\lv_2}\equiv\int\frac{\dd^2\lv_1}{(2\pi)^2}\frac{\dd^2\lv_2}{(2\pi)^2}\;\cD_S(\lv_1)\,\cD_S(\lv_2)\,.
\label{eq:measure}
\end{equation}

For the emission by the antiquark we change variables from $(\x,\y,\z)$ to $(\cc,\R,\Y)$, with unit
Jacobian,
\begin{equation}
\x=\cc+(1-\zeta)\R\,,\qquad\y=\cc-\zeta\R+\frac{\xi}{1-\zeta}\Y\,,\qquad
\z=\cc-\zeta\R-\frac{\eta}{1-\zeta}\Y\,,\qquad\yp=\cc-\zeta\R\,,
\label{eq:var-qbar}
\end{equation}
and for the emission by the quark to $(\cc,\Rt,\X)$,
\begin{equation}
\y=\cc-(1-\eta)\Rt\,,\qquad\x=\cc+\eta\Rt+\frac{\xi}{1-\eta}\X\,,\qquad
\z=\cc+\eta\Rt-\frac{\zeta}{1-\eta}\X\,,\qquad\xp=\cc+\eta\Rt\,.
\label{eq:var-q}
\end{equation}
The phase of the cross-section becomes
\begin{equation}
\kk_1\cdot\x+\kk_2\cdot\y+\kk_3\cdot\z=(\kk_1+\kk_2+\kk_3)\cdot\cc+\kt\cdot\R-\pt\cdot\Y
=(\kk_1+\kk_2+\kk_3)\cdot\cc+\ktq\cdot\Rt-\ptq\cdot\X\,,
\end{equation}
with the relative momenta
\begin{align}
\kt&=(1-\zeta)\kk_1-\zeta(\kk_2+\kk_3)\,,&\pt&=\frac{\eta\kk_3-\xi\kk_2}{1-\zeta}\,,
\nonumber\\
\ktq&=\eta(\kk_1+\kk_3)-(1-\eta)\kk_2\,,&\ptq&=\frac{\zeta\kk_3-\xi\kk_1}{1-\eta}\,.
\label{eq:relmom}
\end{align}
By construction, the effective vertices are the Fourier transforms of momentum-space vertices,
\begin{equation}
\cM^{lr}_{1,2}(\R,\Y)=\int\frac{\dd^2\kk}{2\pi}\frac{\dd^2\pp}{2\pi}\;e^{i\kk\cdot\R-i\pp\cdot\Y}\;
\Phi^{lr}_{1,2}(\kk,\pp)\,,
\end{equation}
where, with the helicity indices $\lambda_1\lambda_2$ suppressed,
\begin{align}
\Phi^{lr}_1(\kk,\pp)&=N^{lr}(\kk,\pp)\,f_1(\kk,\pp)
-\xi\eta\,\big(\delta^{lr}+2i\ee^{lr}\lambda_1\big)\,\delta_{\lambda_1\lambda_2}\;g(\kk,\pp)\,,
\nonumber\\
\Phi^{lr}_2(\kk,\pp)&=N^{lr}(\kk,\pp)\,f_2(\kk,\pp)\,,
\nonumber\\
N^{lr}(\kk,\pp)&\equiv\sum_\lambda\Big[\varphi^{lj}_{\lambda_1\lambda}(\zeta)\,\tau^{rn}_{\lambda\lambda_2}(\xi,\eta)\,k^jp^n
+\xi\of\,\varphi^{lj}_{\lambda_1\lambda}(\zeta)\,\mu^r_{\lambda\lambda_2}\,k^j
\nonumber\\
&\qquad\quad+\mf\,\mu^l_{\lambda_1\lambda}\,\tau^{rn}_{\lambda\lambda_2}(\xi,\eta)\,p^n
+\mf\,\xi\of\,\mu^l_{\lambda_1\lambda}\,\mu^r_{\lambda\lambda_2}\Big]\,,
\label{eq:Phi}
\end{align}
and the scalar functions are the energy denominators,
\begin{align}
f_1(\kk,\pp)&=\frac{1-\zeta}{\big(\kk^2+\ef^2\big)\Big[\zeta(1-\zeta)^2\big(\pp^2+\of^2\big)+\xi\eta\big(\kk^2+\ef^2\big)\Big]}\,,
\nonumber\\
f_2(\kk,\pp)&=\frac{1}{\zeta(1-\zeta)\big(\kk^2+\ef^2\big)\big(\pp^2+\of^2\big)}\,,
\nonumber\\
g(\kk,\pp)&=\frac{1}{\zeta(1-\zeta)^2\big(\pp^2+\of^2\big)+\xi\eta\big(\kk^2+\ef^2\big)}\,,
\nonumber\\
f_1-f_2&=-\frac{\xi\,\eta\;g(\kk,\pp)}{\zeta(1-\zeta)\big(\pp^2+\of^2\big)}\,.
\label{eq:f12}
\end{align}
For the emission by the quark, $\tM_{1,2}(\Rt,\X)$ are the transforms of $\widetilde\Phi_{1,2}$, with
\begin{align}
\widetilde\Phi^{lr}_1(\kk,\pp)&=\widetilde N^{lr}(\kk,\pp)\,\tilde f_1(\kk,\pp)
-\xi\zeta\,\big(\delta^{lr}-2i\ee^{lr}\lambda_1\big)\,\delta_{\lambda_1\lambda_2}\;\tilde g(\kk,\pp)\,,
\nonumber\\
\widetilde\Phi^{lr}_2(\kk,\pp)&=\widetilde N^{lr}(\kk,\pp)\,\tilde f_2(\kk,\pp)\,,
\nonumber\\
\widetilde N^{lr}(\kk,\pp)&\equiv\sum_\lambda\Big[\tau^{nr}_{\lambda_1\lambda}(\xi,\zeta)\,\varphi^{lj}_{\lambda\lambda_2}(\zeta+\xi)\,k^jp^n
+\xi\oft\,\mu^r_{\lambda_1\lambda}\,\varphi^{lj}_{\lambda\lambda_2}(\zeta+\xi)\,k^j
\nonumber\\
&\qquad\quad+\mf\,\tau^{nr}_{\lambda_1\lambda}(\xi,\zeta)\,\mu^l_{\lambda\lambda_2}\,p^n
+\mf\,\xi\oft\,\mu^r_{\lambda_1\lambda}\,\mu^l_{\lambda\lambda_2}\Big]\,.
\label{eq:Phit}
\end{align}
The functions $\tilde f_1$, $\tilde f_2$ and $\tilde g$ are obtained from eq.~\eqref{eq:f12} by
$\zeta\leftrightarrow\eta$, $\ef\to\eft$ and $\of\to\oft$.

\subsection{The amplitudes}

With these definitions the coordinate integrals are elementary. For instance,
\begin{align}
\int\dd^2\R\,\dd^2\Y\;e^{-i\kt\cdot\R+i\pt\cdot\Y}\,\cM_1(\R,\Y)\,S(\x,\z)\,S(\z,\y)
&=(2\pi)^2\int_{\lv_1\lv_2}\Phi_1\big(\kap,\rhov\big)\,,
\nonumber\\
\int\dd^2\R\,\dd^2\Y\;e^{-i\kt\cdot\R+i\pt\cdot\Y}\,\cM_2(\R,\Y)\,S(\x,\yp)
&=(2\pi)^2\int_{\lv_1\lv_2}\Phi_2\big(\kap,\pt\big)\,,
\nonumber\\
\int\dd^2\Rt\,\dd^2\X\;e^{-i\ktq\cdot\Rt+i\ptq\cdot\X}\,\tM_1(\Rt,\X)\,S(\x,\z)\,S(\z,\y)
&=(2\pi)^2\int_{\lv_1\lv_2}\widetilde\Phi_1\big(\kap_q,\rhov_q\big)\,,
\nonumber\\
\int\dd^2\Rt\,\dd^2\X\;e^{-i\ktq\cdot\Rt+i\ptq\cdot\X}\,\tM_2(\Rt,\X)\,S(\xp,\y)
&=(2\pi)^2\int_{\lv_1\lv_2}\widetilde\Phi_2\big(\kap_q,\ptq\big)\,,
\label{eq:FTs}
\end{align}
with the shifted momenta
\begin{equation}
\kap\equiv\kt-\lv_1\,,\quad\;\rhov\equiv\pt+\frac{\eta}{1-\zeta}\,\lv_1-\lv_2\,,\quad\;
\kap_q\equiv\ktq-\lv_2\,,\quad\;\rhov_q\equiv\ptq+\lv_1-\frac{\zeta}{1-\eta}\,\lv_2\,,
\label{eq:shifted-def}
\end{equation}
and where we used the normalisation in eq.~\eqref{eq:DS} to write all the terms under the same double
integral. Here $\lv_1$ is the momentum transferred to the quark and $-\lv_2$ the one transferred to
the antiquark, so that the gluon receives $\lv_2-\lv_1$. The integration over $\cc$ and $\bar\cc$
yields $(2\pi)^2\delta^{(2)}(\kk_1+\kk_2+\kk_3)\,S_\perp$, with $S_\perp$ the transverse area of the
target: for a homogeneous target the coherent process transfers no net transverse momentum. On the
support of this $\delta$-function, the shifted momenta \eqref{eq:shifted-def} are simply the relative
momenta \eqref{eq:relmom} evaluated for the partons before the scattering, i.e.\ with momenta
$\kk_1-\lv_1$, $\kk_2+\lv_2$ and $\kk_3+\lv_1-\lv_2$:
\begin{align}
\kap&=\kk_1-\lv_1\,,&
\rhov&=(\kk_3+\lv_1-\lv_2)+\frac{\xi}{1-\zeta}(\kk_1-\lv_1)\,,
\nonumber\\
\kap_q&=-(\kk_2+\lv_2)\,,&
\rhov_q&=(\kk_3+\lv_1-\lv_2)+\frac{\xi}{1-\eta}(\kk_2+\lv_2)\,.
\label{eq:shifted}
\end{align}
The coherent amplitude for the emission by the antiquark is thus built from the following five
transverse structures, one for each helicity path and one for the instantaneous term,
\begin{align}
T^{jn}&=\int_{\lv_1\lv_2}\Big[\kappa^j\rho^n\,f_1(\kap,\rhov)-\kappa^j\,\tilde p^{\,n}\,f_2(\kap,\pt)\Big]
-\tilde k^j\tilde p^{\,n}\Big[f_1(\kt,\pt)-f_2(\kt,\pt)\Big]\,,
\nonumber\\
T_k^{j}&=\int_{\lv_1\lv_2}\Big[\kappa^j\,f_1(\kap,\rhov)-\kappa^j\,f_2(\kap,\pt)\Big]
-\tilde k^j\Big[f_1(\kt,\pt)-f_2(\kt,\pt)\Big]\,,
\nonumber\\
T_p^{n}&=\int_{\lv_1\lv_2}\Big[\rho^n\,f_1(\kap,\rhov)-\tilde p^{\,n}\,f_2(\kap,\pt)\Big]
-\tilde p^{\,n}\Big[f_1(\kt,\pt)-f_2(\kt,\pt)\Big]\,,
\nonumber\\
T_0&=\int_{\lv_1\lv_2}\Big[f_1(\kap,\rhov)-f_2(\kap,\pt)\Big]-\Big[f_1(\kt,\pt)-f_2(\kt,\pt)\Big]\,,
\qquad
T_I=\int_{\lv_1\lv_2}g(\kap,\rhov)-g(\kt,\pt)\,.
\label{eq:T}
\end{align}
In each of them the first term is the emission before the collision followed by the scattering of the
three partons, the second term the scattering of the pair followed by the emission, and the last term
the no-scattering subtraction; all of them vanish when $\cD_S(\lv)\to(2\pi)^2\delta^{(2)}(\lv)$, that is,
in the absence of scattering. The structures $\widetilde T^{jn}$, $\widetilde T_k^j$,
$\widetilde T_p^n$, $\widetilde T_0$ and $\widetilde T_I$ for the emission by the quark are given by the
same formulas with $(\kap,\rhov)\to(\kap_q,\rhov_q)$, $(\kt,\pt)\to(\ktq,\ptq)$ and
$(f_1,f_2,g)\to(\tilde f_1,\tilde f_2,\tilde g)$. In terms of these structures the amplitudes read
\begin{align}
\Amp^{lr}_{\lambda_1\lambda_2}=\sum_\lambda\Big[&\varphi^{lj}_{\lambda_1\lambda}(\zeta)\,\tau^{rn}_{\lambda\lambda_2}(\xi,\eta)\,T^{jn}
+\xi\of\,\varphi^{lj}_{\lambda_1\lambda}(\zeta)\,\mu^r_{\lambda\lambda_2}\,T_k^j
+\mf\,\mu^l_{\lambda_1\lambda}\,\tau^{rn}_{\lambda\lambda_2}(\xi,\eta)\,T_p^n
\nonumber\\
&+\mf\,\xi\of\,\mu^l_{\lambda_1\lambda}\,\mu^r_{\lambda\lambda_2}\,T_0\Big]
-\xi\eta\,\big(\delta^{lr}+2i\ee^{lr}\lambda_1\big)\,\delta_{\lambda_1\lambda_2}\,T_I\,,
\nonumber\\
\widetilde\Amp^{lr}_{\lambda_1\lambda_2}=\sum_\lambda\Big[&\tau^{nr}_{\lambda_1\lambda}(\xi,\zeta)\,\varphi^{lj}_{\lambda\lambda_2}(\zeta+\xi)\,\widetilde T^{jn}
+\xi\oft\,\mu^r_{\lambda_1\lambda}\,\varphi^{lj}_{\lambda\lambda_2}(\zeta+\xi)\,\widetilde T_k^j
+\mf\,\tau^{nr}_{\lambda_1\lambda}(\xi,\zeta)\,\mu^l_{\lambda\lambda_2}\,\widetilde T_p^n
\nonumber\\
&+\mf\,\xi\oft\,\mu^r_{\lambda_1\lambda}\,\mu^l_{\lambda\lambda_2}\,\widetilde T_0\Big]
-\xi\zeta\,\big(\delta^{lr}-2i\ee^{lr}\lambda_1\big)\,\delta_{\lambda_1\lambda_2}\,\widetilde T_I\,,
\label{eq:amp}
\end{align}
and the coordinate integrals of the elastic part of the cross-section give
\begin{align}
\int_{\x,\y,\z,\bx,\by,\bz}&e^{-i\kk_1\cdot(\x-\bx)-i\kk_2\cdot(\y-\by)-i\kk_3\cdot(\z-\bz)}\,\Big[\ldots\Big]_{\rm el}
\nonumber\\
&=S_\perp\,(2\pi)^2\delta^{(2)}(\kk_1+\kk_2+\kk_3)\;(2\pi)^4\sum_{\lambda_1,\lambda_2}\sum_{l,r}
\Amp^{lr}_{\lambda_1\lambda_2}\,\Amp^{lr\,*}_{\lambda_1\lambda_2}\,,
\label{eq:coordint}
\end{align}
for the antiquark, and similarly for the other two channels.

\subsection{Results}

The helicity and polarisation sums are the momentum-space counterparts of
eqs.~\eqref{eq:H-direct}--\eqref{eq:Hx}. We use the notations
$|T|^2\equiv\sum_{j,n}T^{jn}T^{jn\,*}$, $\tr\,T\equiv T^{11}+T^{22}$,
$T_k\cdot T_p^*\equiv\sum_jT_k^jT_p^{j\,*}$ and
\begin{equation}
\ee^{ij}\ee^{mn}\,T^{im}U^{jn\,*}=T^{11}U^{22\,*}+T^{22}U^{11\,*}-T^{12}U^{21\,*}-T^{21}U^{12\,*}\,.
\end{equation}

\paragraph{Antiquark.} With $\sum\Amp\Amp^*=H^{\rm massless}+H^{\rm massive}$,
\begin{align}
H^{\rm massless}={}&8\big[\zeta^2+(1-\zeta)^2\big]\big[\eta^2+(1-\zeta)^2\big]\,|T|^2
-8\,\xi\,(2\zeta-1)\big(2(1-\zeta)-\xi\big)\,\ee^{ij}\ee^{mn}\,T^{im}T^{jn\,*}
\nonumber\\
&-16\,\zeta\,\xi\,\eta^2\;\re\big[\tr\,T\;T_I^*\big]+8\,\xi^2\eta^2\,|T_I|^2\,,
\nonumber\\
H^{\rm massive}={}&8\,\mf^2\big[\eta^2+(1-\zeta)^2\big]\,T_p\cdot T_p^*
+8\,\xi^2\of^2\big[\zeta^2+(1-\zeta)^2\big]\,T_k\cdot T_k^*+8\,\mf^2\xi^2\of^2\,|T_0|^2
\nonumber\\
&-16\,\zeta\,\eta\,\mf\,\xi\of\;\re\big[\tr\,T\;T_0^*-T_k\cdot T_p^*\big]
+16\,\xi^2\eta\,\mf\,\of\;\re\big[T_0\,T_I^*\big]\,,
\label{eq:H-el-qbar}
\end{align}
the coherent contribution to the cross-section reads
\begin{align}
\frac{\dd\sigma^{(\bar q\bar q)}_{\rm el}}{\dd^3k_1\,\dd^3k_2\,\dd^3k_3}={}&
\frac{\aem\,\as\,\CF\Nc\,S_\perp}{2(2\pi)^{4}}\;\frac{\zeta^2}{\xi\,(q^+)^2}\;
\delta\big(q^+-k_1^+-k_2^+-k_3^+\big)\,\delta^{(2)}(\kk_1+\kk_2+\kk_3)
\nonumber\\
&\times\sum_fe_f^2\;\Big[H^{\rm massless}+H^{\rm massive}\Big]\,.
\label{eq:xs-el-qbar}
\end{align}

\paragraph{Quark.} With $\sum\widetilde\Amp\widetilde\Amp^*=\widetilde H^{\rm massless}+\widetilde H^{\rm massive}$,
\begin{align}
\widetilde H^{\rm massless}={}&8\big[\zeta^2+(1-\eta)^2\big]\big[\eta^2+(1-\eta)^2\big]\,|\widetilde T|^2
-8\,\xi\,(2\eta-1)\big(2(1-\eta)-\xi\big)\,\ee^{ij}\ee^{mn}\,\widetilde T^{im}\widetilde T^{jn\,*}
\nonumber\\
&+16\,\zeta^2\xi\,\eta\;\re\big[\tr\,\widetilde T\;\widetilde T_I^*\big]+8\,\xi^2\zeta^2\,|\widetilde T_I|^2\,,
\nonumber\\
\widetilde H^{\rm massive}={}&8\,\mf^2\big[\zeta^2+(1-\eta)^2\big]\,\widetilde T_p\cdot\widetilde T_p^*
+8\,\xi^2\oft^{\,2}\big[\eta^2+(1-\eta)^2\big]\,\widetilde T_k\cdot\widetilde T_k^*
+8\,\mf^2\xi^2\oft^{\,2}\,|\widetilde T_0|^2
\nonumber\\
&+16\,\zeta\,\eta\,\mf\,\xi\oft\;\re\big[\tr\,\widetilde T\;\widetilde T_0^*-\widetilde T_k\cdot\widetilde T_p^*\big]
+16\,\xi^2\zeta\,\mf\,\oft\;\re\big[\widetilde T_0\,\widetilde T_I^*\big]\,,
\label{eq:H-el-q}
\end{align}
\begin{align}
\frac{\dd\sigma^{(qq)}_{\rm el}}{\dd^3k_1\,\dd^3k_2\,\dd^3k_3}={}&
\frac{\aem\,\as\,\CF\Nc\,S_\perp}{2(2\pi)^{4}}\;\frac{\eta^2}{\xi\,(q^+)^2}\;
\delta\big(q^+-k_1^+-k_2^+-k_3^+\big)\,\delta^{(2)}(\kk_1+\kk_2+\kk_3)
\nonumber\\
&\times\sum_fe_f^2\;\Big[\widetilde H^{\rm massless}+\widetilde H^{\rm massive}\Big]\,.
\label{eq:xs-el-q}
\end{align}

\paragraph{Interference.} With $\sum\widetilde\Amp\Amp^*=H_\times^{\rm massless}+H_\times^{\rm massive}$,
\begin{align}
H_\times^{\rm massless}={}&8\big[1-\xi-2\zeta\eta\big]\big[\xi(1-\xi)+2\zeta\eta\big]\sum_{j,n}\widetilde T^{jn}T^{jn\,*}
-8\,\xi\,\big[1-2\zeta-\xi\big]^2\,\ee^{ij}\ee^{mn}\,\widetilde T^{im}T^{jn\,*}
\nonumber\\
&-8\,\xi\,\eta\,(\zeta+\xi)^2\;\tr\,\widetilde T\;T_I^*+8\,\xi\,\zeta\,(1-\zeta)^2\;\widetilde T_I\;\tr\,T^*\,,
\nonumber\\
H_\times^{\rm massive}={}&8\,\mf^2\,(\xi+2\zeta\eta)\,\widetilde T_p\cdot T_p^*
-8\,(\xi+2\zeta\eta)\,\xi^2\,\oft\,\of\;\widetilde T_k\cdot T_k^*
\nonumber\\
&-8\,(\zeta+\xi)^2\,\mf\,\xi\of\;\tr\,\widetilde T\;T_0^*+8\,(1-\zeta)^2\,\mf\,\xi\oft\;\widetilde T_0\;\tr\,T^*
\nonumber\\
&-8\,\eta^2\,\mf\,\xi\oft\;\widetilde T_k\cdot T_p^*+8\,\zeta^2\,\mf\,\xi\of\;\widetilde T_p\cdot T_k^*\,,
\label{eq:H-el-int}
\end{align}
\begin{align}
\frac{\dd\sigma^{(q\bar q)}_{\rm el}}{\dd^3k_1\,\dd^3k_2\,\dd^3k_3}={}&-
\frac{\aem\,\as\,\CF\Nc\,S_\perp}{2(2\pi)^{4}}\;\frac{\zeta\,\eta}{\xi\,(q^+)^2}\;
\delta\big(q^+-k_1^+-k_2^+-k_3^+\big)\,\delta^{(2)}(\kk_1+\kk_2+\kk_3)
\nonumber\\
&\times\sum_fe_f^2\;2\,\re\Big[H_\times^{\rm massless}+H_\times^{\rm massive}\Big]\,.
\label{eq:xs-el-int}
\end{align}
The labels ``massless'' and ``massive'' refer to the tensor structures: the former are those which are
already present for $\mf=0$, but they still depend on the mass through the scales $\ef$, $\of$ and
their tilded counterparts in the functions $f_1$, $f_2$ and $g$. Equations
\eqref{eq:xs-el-qbar}--\eqref{eq:xs-el-int} reduce the coherent trijet cross-section to two
two-dimensional integrals over the momenta transferred by the target, eq.~\eqref{eq:T}, which are well
suited for a numerical evaluation with any model for the dipole $S$-matrix.
